\documentclass[11pt]{amsart}
\usepackage{geometry}                % See geometry.pdf to learn the layout options. There are lots.
\geometry{letterpaper}                   % ... or a4paper or a5paper or ... 
%\geometry{landscape}                % Activate for for rotated page geometry
%\usepackage[parfill]{parskip}    % Activate to begin paragraphs with an empty line rather than an indent
\usepackage{graphicx}
\usepackage{amssymb}
\usepackage{epstopdf}
\DeclareGraphicsRule{.tif}{png}{.png}{`convert #1 `dirname #1`/`basename #1 .tif`.png}

\title{Problema 3.15}
\author{}
%\date{}                                           % Activate to display a given date or no date

\begin{document}
\maketitle
%\section{}
%\subsection{}
Haga un diagrama de flujo para calcular lo que hay que pagar por un conjunto de llamadas telefónicas. Por cada llamada se ingresa el tipo (Internacional, Nacional, Local) y la duración en minutos. El criterio que se sigue para calcular el costo de cada llamada es el siguiente: \\

\indent Internacional: 3 primeros minutos \$7.59 Cada minuto adicional \$3.03 \\

\indent Nacional: 3 primeros minutos \$1.20 Cada minuto adicional \$0.48 \\

\indent Local: Las primeras 50 llamadas no se cobran. Luego, cada llamada cuesta \$0.60 \\

Datos: TIPOi, DURi, TIP02, DUR2,...,X,-1 \\ 

Donde: \\
	\indent TIPO¡ es una variable de tipo caracter que expresa el tipo de la llamada i. 
     Toma el valor de ‘I’ si la llamada es internacional, ‘N’ si es nacional y ‘L si es local. \\

	\indent DURi es una variable de tipo entero que indica la duración de la llamada i en minutos.\\
	
\textbf{Explicación de las variables}\\
CL: Variable de tipo entero. Acumula el número de llamadas locales.\\
CUENTA: Variable de tipo real. Acumula el costo de cada llamada.\\
TIPO: Variable de tipo caracter.\\
DUR: Variable de tipo entero.\\
COSTO: Variable de tipo real. Almacena el costo de cada llamada.\\
En la siguiente tabla se puede observar el seguimiento del algoritmo para el
siguiente conjunto de llamadas.\\
Datos: I, 7, N, 6, N, 12 ,L ,5 ,L ,7 ,L ,15, N, 16, L, 7, L, 6, L, 4 ,1, 11, X, -1\\  


\begin{center}
\begin{tabular}{|c|c|c|c|c|}
\hline
\multicolumn{5}{ |c| }{Tabla 3.17} \\ \hline
  TIPO & DUR & CL & COSTO & CUENTA \\ \hline
  
  I & 7 & 0 & & 0 \\\hline
  
  N & 6 &  & 19.71 & 19.71 \\ \hline
  
  N & 12 &  & 2.64 & 022.35 \\ \hline
  
  L & 5 &  & 5.52 & 27.87 \\ \hline
  
  L & 7 & 1 & 0 & 27.87\\ \hline
  
  L & 15 & 2 & 0 & 27.87 \\ \hline
  
  N & 16 & 3 & 0 & 27.87 \\ \hline
  
  L & 7 &  & 7.44 & 35.31 \\ \hline
  
  L & 6 & 4 & 0 & 35.31 \\ \hline
  
  L & 4 & 5 & 0 & 35.31 \\ \hline
  
  I & 11 & 6 & 0 & 35.31 \\ \hline
  
  X & -1 &  & 31.83 & \underline {67.14} \\ \hline
\end{tabular}
\end{center}

\end{document}  